\chapter{Introduction}\label{chap:intro}

\section{Motivation}\label{sec:motivation}

In a small classroom, a good teacher rarely relies on test scores alone. A furrowed brow during a derivation, a distant gaze halfway through an explanation, a brief smile after a difficult problem, and similar cues quietly tell the teacher when to slow down, offer a hint, adjust the pace, switch explanatory strategy, or offer a word of encouragement. Decades of research in affective neuroscience and educational psychology confirm that this teacherly intuition has deep roots. Emotion is not a byproduct of thinking but an integral part of it, and learning, as it unfolds in schools and in the real world, is not a disembodied or purely rational process \cite{immordino-yang_we_2007}. Emotional states modulate attention, encoding, and memory consolidation, so that what a student feels while studying directly shapes what they later remember and how they will reason about it \cite{tyng_influences_2017}. In tutoring contexts specifically, a distinctive family of academic emotions, namely engagement, confusion, frustration, and boredom, emerges repeatedly, each carrying pedagogical meaning of its own \cite{dmello_language_2012}. Pekrun's Control-Value Theory explains \emph{why} these emotions matter: fluctuations in learners' appraisals of control over a task and of its personal value arouse achievement emotions that alter how information is processed, stored, and applied, feeding back into engagement, motivation, and ultimately performance \cite{pekrun_control-value_2006}. A learner's face, therefore, is not decorative information; it is a window into the very processes that education is trying to shape.

However, the contextual cues that a skilled teacher reads almost unconsciously largely disappear when learning moves online or to large-scale digital platforms. Modern e-learning and intelligent tutoring environments have expanded rapidly, especially after the pandemic, and are now a widely used channel for student learning\cite{strielkowski_span_2025}. Yet these systems typically treat learners as receivers of information, with little concern for their moment-to-moment emotional or cognitive condition. Under Pekrun's Control-Value Theory, this neglect carries direct cognitive costs: whilst positive emotions such as enjoyment sustain motivation and broaden attentional resources, negative states such as anxiety and boredom produce cognitive interference that impairs how information is encoded and retained \cite{pekrun_control-value_2006}. The academic emotions most characteristic of tutoring interactions, namely confusion, frustration, and boredom, identified by D'Mello and Graesser as the states that recur most reliably in learning dialogues \cite{dmello_language_2012}, can therefore persist silently and significantly hinder learning performance \cite{gong_design_2025}. Studies of online and STEM learning repeatedly show that emotions shape engagement, motivation, and satisfaction, and that ignoring them leaves platforms unable to respond when students disengage \cite{anwar_emotions_2023}. Effective tutoring responds to such shifts not merely by selecting a harder or easier item, but by adjusting pace, offering scaffolding or hints, switching explanatory strategy, or providing motivational support; this multi-dimensional responsiveness is precisely what digital platforms most lack. The effect is compounded by the fact that supervision is thin in online settings: when nobody notices that attention is drifting, sustained inattention tends to be one of the first things that breaks down, eroding the very benefits that the medium was meant to deliver \cite{gupta_edfa_2023,aly_enhancing_2023}. A system capable of reading the learner's face in real time would offer gains on several of these fronts simultaneously: it could detect confusion or boredom well before a failing assessment reveals the problem; adjust content difficulty, pacing, and support strategies continuously for each individual rather than for the average learner; assist teachers in monitoring emotional engagement across large cohorts at once; and generate richer behavioural and affective data than click-logs or response times alone can provide. The result would more closely resemble the personalised quality of one-to-one tutoring, and it would align with what learning science recommends: timely corrective feedback, guided reflection, and conditions that support long-term retention. The picture that emerges is therefore a paradox: precisely in the environments where such personalisation across content, pacing, and learner support is most needed, the richest source of human feedback, the learner's face, is the one that digital systems most systematically ignore.

\section{Problem Statement}\label{sec:problem}

A widely proposed response to this paradox is adaptive learning: systems that tailor the sequence, difficulty, feedback, and pacing of material to each individual learner rather than delivering a fixed path \cite{gligorea_adaptive_2023,strielkowski_span_2025}. Recent \ac{ITS} have taken up this promise with increasingly sophisticated tools. Reinforcement-learning approaches have been used to adjust task difficulty and instructional complexity in simulated learners \cite{axak_adaptive_2025,deshmukh_developing_2025}, and have been combined with knowledge tracing to optimise personalised learning paths across students' evolving states \cite{fu_integrating_2025}. Comparative studies show that such policies can meaningfully improve task selection over rule-based baselines \cite{azoulay_adaptive_2021}, and recent work couples knowledge tracing with reinforcement-learning-based problem selection to produce measurable learning gains in short human studies \cite{guzman_developing_2025}. Even outside the reinforcement-learning paradigm, biologically inspired memory architectures have been used to recommend the next exercise within each learner's Zone of Proximal Development \cite{kuling_ken_2025}. Despite this progress, two patterns emerge clearly from recent surveys of the field. First, a surprising number of adaptive systems are described in the literature but remain experimental, with many not validated in authentic classroom settings and often disconnected from the concrete problems students face \cite{kabudi_ai-enabled_2021,zerkouk_comprehensive_2025}. Second, and more importantly for this thesis, most adaptive systems still operate only on cognitive-behavioural traces, such as answers, response times, clicks, and hint requests, and treat the learner as affect-free \cite{axak_adaptive_2025,fu_integrating_2025,azoulay_adaptive_2021}. Broader comparative work echoes that behavioural traces predominate when signals are chosen for the policy loop \cite{guzman_developing_2025,weitekamp_tutorgym_2025}. Urbaite warns that a persistent risk in such systems is collapsing the multidimensional construct of engagement into ``easy-to-log proxies'' such as clicks and time-on-page \cite{urbaite_adaptive_2026}. Meng and Yang similarly observe that existing \ac{RL}-driven tutoring solutions often overlook engagement factors altogether and, when they do attempt to model them, rely primarily on textual interaction patterns rather than on perceptual signals of the student \cite{meng_self-evolving_2025}. A related but distinct limitation is that the action spaces of most \ac{RL}-based tutors remain narrow by design: policies typically select among difficulty levels or choose which item to present next, but do not cover the broader pedagogical repertoire that a skilled tutor employs, including scaffolding, hints, pacing changes, motivational messages, and explanatory reformulations. Authors of recent \ac{RL}-based tutors explicitly flag the integration of affective computing as open future work, acknowledging that their policies are, for now, both emotionally blind and pedagogically narrow \cite{axak_adaptive_2025,guzman_developing_2025}.

\ac{FER} offers a practical way to recover part of the signal that adaptive systems currently ignore. Over the last decade, deep learning has substantially advanced \ac{FER} accuracy, with convolutional networks, vision transformers, and hybrid temporal models all contributing to stronger recognition in unconstrained settings \cite{lek_academic_2023,liu_dynamic_2026,aly_revolutionizing_2025}. Attention-heavy variants further tighten performance in challenging conditions \cite{aly_comprehensive_2025}. In educational contexts, researchers have moved beyond Ekman's six basic emotions, which rarely occur in actual study sessions, towards academic affect dimensions such as engagement, boredom, confusion, and frustration, the very states singled out by earlier discourse-analytic work on tutoring \cite{leong_deep_2020,dmello_language_2012}. The \ac{DAiSEE} dataset has become a common anchor for this line of research, and a growing body of deep-learning models, ranging from facial-embedding and landmark approaches to more recent dynamic \ac{FER} models with attention and temporal mechanisms, has been trained on it \cite{leong_deep_2020,asaju_adaptive_2022,mehta_three-dimensional_2022}. Related architectures and training regimes continue to expand in the latest \ac{DAiSEE}-based studies \cite{komaravalli_detecting_2023,liu_dynamic_2026,rao_recognition_2021}. Parallel efforts in online-learning contexts use modern \ac{CNN} and transformer architectures to monitor students in live classes \cite{aly_revolutionizing_2025,aly_enhancing_2023}. Additional classroom-oriented deployments extend the same design space \cite{aly_comprehensive_2025,gupta_edfa_2023}. Systematic reviews of academic \ac{FER} confirm that deep learning now dominates this subfield and that real-time deployment is becoming realistic, while also noting ongoing challenges from pose, illumination, occlusion, and cross-cultural variability \cite{lek_academic_2023,fernandez-herrero_evaluating_2024}. Overall, the evidence that facial signals carry meaningful information about engagement and cognitive state in online learners has become difficult to dispute \cite{mehta_three-dimensional_2022,gupta_edfa_2023,woolf_face_2023}; yet accurate \ac{FER} is necessary but insufficient for adaptive tutoring unless those signals are coupled to concrete instructional decisions that extend beyond difficulty adjustment or next-item selection alone.

Some recent work does connect affect to adaptation, including reinforcement-learning formulations that map inferred emotional state to instructional adjustments such as difficulty, pacing, scaffolding, hints, or metacognitive prompts. In practice, however, many of these systems lean on multimodal or physiological channels, such as heart rate, EEG, respiration, text, voice, or tone \cite{gong_design_2025,govea_implementation_2024}. Such setups are often heavy in practice, since signals must be acquired, synchronised, and processed alongside the teaching interaction, which typically implies specialised hardware and careful integration before the approach can be deployed at scale in ordinary online courses \cite{gutierrez_development_2025,kong_real-time_2025,nandimandalam_ai-powered_2025}. A second limitation, within this same emotion-aware adaptation literature, is that affect is not always carried through to the decision process that selects a pedagogical intervention; many systems stop at detection or dashboard-style reporting, so the intended loop from detecting emotion, to using it in the adaptation decision, to executing a concrete instructional action is rarely closed \cite{fernandez-herrero_evaluating_2024,septiana_emotion-related_2024,vistorte_integrating_2024}. Using facial expression alone offers a lighter alternative that is less intrusive than body-worn sensors and realistic to collect with commodity webcams for large cohorts. Prior reinforcement-learning work that uses facial affect as feedback for adaptation shows that such a loop is feasible, but evaluation is often simulation-based, the affective signal may be coarse or averaged across whole tasks rather than supplied as a continuous input at decision time, and the policy action space typically remains narrow \cite{perez_emotions_2024}. This thesis therefore adopts an \ac{FER}-only, webcam-based pipeline in which facial affect drives real-time pedagogical action selection through a reinforcement-learning policy over a multi-action instructional space, moving beyond the difficulty-centric sequencing envisaged in earlier project iterations.

While facial expression recognition in education and adaptive tutoring have both matured, they have largely matured in parallel. Most \ac{FER}-for-learning systems stop at recognition: they train a classifier on \ac{DAiSEE} or a similar dataset, report per-class accuracy, and at best propose an analytics dashboard or teacher-facing feedback channel \cite{leong_deep_2020,rao_recognition_2021,mehta_three-dimensional_2022}. Studies that span \ac{DAiSEE}-trained models to online prototypes typically remain at the recognition stage for design reasons \cite{komaravalli_detecting_2023,liu_dynamic_2026,aly_revolutionizing_2025}. Further deployments discussed in the same literature often mirror that scope \cite{aly_enhancing_2023}. Where these systems do discuss adaptation, it is typically aspirational, leaving the pedagogical response either to a human instructor or to unspecified future work \cite{asaju_adaptive_2022,aly_comprehensive_2025,Prabakaran2025FacialExpressionAdaptiveTeaching}. In parallel, adaptive tutoring systems that do implement closed-loop decisions rarely use facial affect on its own; as noted above, the face, when included, tends to be one channel among several, with non-facial signals typically carrying most of the weight. Scoping reviews make the same observation at a field level. Affective intelligent tutoring systems, it has been noted, ``recognize more than they regulate'', still leaning on single-modality detection and rarely closing the loop with real-time interventions \cite{yuvaraj_affective_2025}. Much of the underlying emotion-recognition work, moreover, still relies on static datasets and predefined viewing scenarios, which can underrepresent how academic affect fluctuates during live online study. There remains, in short, a lack of complete pipelines in which webcam-based academic-affect recognition drives a closed-loop instructional policy: a reinforcement-learning tutor that selects among diverse pedagogical interventions rather than limiting adaptation to the difficulty of the next exercise alone, together with an evaluation design that isolates the contribution of the affective channel. This thesis is motivated by that gap.
\section{Objectives}\label{sec:objectives}

This thesis develops an adaptive intelligent tutoring system that couples real-time recognition of learner affect with a sequential decision-making policy that selects pedagogical interventions in response to the learner's state. The objectives below state, for each component and for the integrated system, what the work set out to achieve and how its success is measured.

\begin{enumerate}
  \item To build and evaluate a standalone multi-label academic \ac{FER} model on \ac{DAiSEE} that infers the co-occurring affective states of boredom, engagement, confusion, and frustration from video of the learner \cite{gupta_daisee_2016,mehta_three-dimensional_2022}. The model performs multi-label affective state inference with particular emphasis on the reliable detection of minority affective states under the severe class imbalance of the benchmark \cite{aly_comprehensive_2025}. Recognition performance is evaluated under a single standardised evaluation protocol using per-label and aggregate F1, accuracy, precision, and recall, so that the reported results reflect model quality rather than differences in data preparation.
  \item To formulate adaptive tutoring as a discrete-time \ac{MDP} that represents the learner's evolving cognitive and affective state and selects among a broad repertoire of pedagogical interventions rather than difficulty adjustment alone \cite{axak_adaptive_2025,azoulay_adaptive_2021}. The transition dynamics and the multi-objective reward are grounded in established cognitive and affective learning theory, with the reward jointly optimising cognitive learning gain and affective well-being. Success is assessed by whether the formulation yields a well-posed, theory-consistent decision problem capable of supporting policy learning.
  \item To learn and evaluate a value-based reinforcement-learning tutoring policy that selects pedagogical interventions from the learner's observed state so as to maximise the multi-objective reward \cite{mnih_dqn_2015}. The learned policy is evaluated in simulation against an expert rule-based tutor and a random baseline on normalised learning gain, episode success rate, cumulative reward, consistency with theory-derived emotion-conditioned interventions, and the diversity of its action selection \cite{fernandez-herrero_evaluating_2024,yuvaraj_affective_2025,lek_academic_2023}.
    \item To develop a unified closed-loop architecture integrating facial expression recognition and reinforcement learning, and to evaluate its affect-driven decision-making in relation to established educational and affective theories. \cite{yuvaraj_affective_2025}.
\end{enumerate}
\section{Thesis Outline}\label{sec:outline}

The remainder of this thesis is organised as follows:

\begin{itemize}
  \item \textbf{Background} reviews the theoretical foundations of emotion in learning, facial expression recognition, and adaptive learning methods.
  \item \textbf{Literature Review} surveys prior work on academic emotion detection, adaptive tutoring, and emotion-aware adaptation in intelligent tutoring systems.
  \item \textbf{Methodology} presents the design of the \ac{FER} module, the reinforcement-learning adaptation module (including the Markov decision process formulation, reward function, and learned policies), and their integration into a closed-loop pipeline.
  \item \textbf{Results} reports the experimental setup and quantitative findings, first for the \ac{FER} component in isolation and then for reinforcement-learning policy comparison against expert rule and random baselines.
  \item \textbf{Conclusion} interprets the findings in relation to the gap identified above and discusses limitations and ethical considerations.
  \item \textbf{Future Work} outlines directions for extending the \ac{FER} module and the tutoring system beyond the scope of this thesis.
\end{itemize}
